The first phase of this research establishes the empirical foundation upon which the subsequent artifact design rests. By systematically mapping the translational landscape of AI in geriatric nutrition, this scoping review identifies the specific gaps, limitations, and opportunities that define the problem space --- and, critically, derives the design requirements that the Snap and Say artifact must satisfy.

\section{Introduction}

While recent reviews have synthesised the general landscape of AI in nutrition \citep{theodorearmandApplicationsArtificialIntelligence2024} or highlighted the potential of AI specifically for healthy aging \citep{grootswagersArtificialIntelligenceNutrition2025}, these syntheses often conflate theoretical model development with practical application. There is limited evidence distinguishing purely algorithmic experiments from systems that have been operationally tested or deployed specifically within geriatric settings.

To address this gap, the primary objective of this scoping review is to identify and map the existing translational evidence on the application of deployed or operationally-tested AI systems in nutrition for aging populations. Unlike previous reviews, we focus specifically on systems evaluated in settings involving adults aged 60 and over to characterise the state of applied AI, identifying specific modalities facilitating clinical integration and delineating barriers to adoption. This review addresses the following research questions:

\begin{enumerate}
    \item What are the characteristics of AI systems in geriatric nutrition, including their intended functions, underlying technologies, and stage of deployment?
    \item What data modalities and computational features are used to develop these systems?
    \item What is the availability status of the datasets used for these systems?
    \item In what settings and populations have these AI systems been studied?
    \item How have these systems been evaluated, and what are the reported findings on their performance?
    \item What are the reported challenges and limitations associated with the application of AI in geriatric nutrition?
    \item What future directions are identified in the literature for advancing the field?
\end{enumerate}

A scoping review methodology was selected as it is designed to map the extent, range, and nature of evidence in an emerging field \citep{arkseyScopingStudiesMethodological2005,levacScopingStudiesAdvancing2010,triccoPRISMAExtensionScoping2018}. The review's research questions are exploratory, aiming to clarify concepts and chart the landscape of the field. This contrasts with a systematic review, which typically answers a focused question about intervention effectiveness \citep{arkseyScopingStudiesMethodological2005}.

\section{Methods}

This scoping review was designed and reported in accordance with the Preferred Reporting Items for Systematic reviews and Meta-Analysis extension for Scoping Reviews (PRISMA-ScR) checklist \citep{pollockRecommendationsExtractionAnalysis2023}.

The protocol was developed a priori. It was registered with the Open Science Framework (OSF) on October 2, 2025, and can be accessed at \url{https://osf.io/c9pg5}. No amendments to the protocol have been made since its registration. All supplementary files, including the full data extraction table, are available at this registration.

Eligibility criteria were established using the Population-Concept-Context (PCC) framework, as recommended for scoping reviews \citep{triccoPRISMAExtensionScoping2018}.

\begin{itemize}
    \item \textbf{Population:} Studies focusing on older adults, defined as individuals aged 60 years and over. This threshold was selected to align with definitions from the World Health Organization (WHO) \citep{worldhealthorganizationWorldReportAgeing2015}. Studies with mixed-age populations were included only if data for older adults were reported separately or constituted the majority ($>75\%$) of the sample.
    \item \textbf{Concept:} Application of AI in a nutrition-related context. Eligible studies must describe an evaluated system, tool, or application. Studies that were purely theoretical, conceptual, or described only model development without integration into a user-facing system were excluded. AI was broadly defined to include machine learning, deep learning, natural language processing, computer vision, expert systems, and predictive analytics.
    \item \textbf{Context:} Any geographical or healthcare setting, including community, acute care, long-term care (LTC), and rehabilitation facilities.
\end{itemize}

Peer-reviewed journal articles and full-text conference proceedings published in English between January 2015 and October 2025 were included.

A search was conducted across MEDLINE (via PubMed), Scopus, Web of Science, IEEE Xplore, and the ACM Digital Library. The final search was conducted on October 2, 2025. The search strategy combined Medical Subject Headings (MeSH) and free-text keywords for four core concepts: (1) aging population, (2) nutrition, (3) AI method, and (4) AI system. The full search strategy for MEDLINE (PubMed) is available in Appendix~\ref{app:search-strategy}.

To enhance comprehensiveness, Google Gemini 2.5 Pro was utilised to generate synonyms, alternative phrases, and subject headings. The research team reviewed, validated, and refined AI-generated suggestions before their incorporation into the final search strings.

Unique records were exported to ASReview systematic review software \citep{vandeschootOpenSourceMachine2021} for screening. A single reviewer (F.W.N.S.) conducted a two-stage screening process. Two verification steps were integrated: (1) sources of uncertain eligibility were discussed with the supervisory team (Y.B., B.T.) to reach consensus, and (2) a second reviewer (B.T.) independently assessed a random $20\%$ sample of excluded full-text articles.

A data charting form, developed in Google Forms, was piloted on five studies and refined with feedback from the supervisory team. The primary reviewer charted all data. For quality assurance, a second reviewer (Y.B.) independently verified a random $20\%$ subset of the charted data against the full-text articles.

The following data items were extracted: publication details, study characteristics, population characteristics, AI details, system and data characteristics (input data modality, dataset availability, computational features, system stage of deployment), nutrition application, evaluation and findings, identified challenges or limitations, and proposed solutions or future directions.

A meta-analysis was not conducted due to heterogeneity. Instead, a narrative synthesis of the charted data was conducted, structured thematically by nutritional application and AI technology.

\section{Results}
\label{sec:results}

The selection process filtered 587 initial records down to 25 included studies based on strict age and empirical evaluation criteria, as illustrated in Figure~\ref{fig:prisma-flow}. Appendix~\ref{app:included-studies} provides a comprehensive characterisation of these 25 systems.

\begin{figure}[htbp]
    \centering
    \includegraphics[width=0.85\textwidth]{1_prisma.pdf}
    \caption{PRISMA 2020 flow diagram detailing the study selection process. Generated using the PRISMA2020 R package.}
    \label{fig:prisma-flow}
\end{figure}

The publication timeline of the 25 included studies (2016--2025) shows that publication frequency increased over time, with $68\%$ ($n=17$) published since 2021, suggesting growing research interest (Figure~\ref{fig:publication-trends}). The studies originated from 18 countries across five continents, primarily Asia ($n=13$, $52\%$) and Europe ($n=9$, $36\%$). Research was concentrated in high-income and upper-middle-income countries, with a notable gap from lower-income settings.

\begin{figure}[htbp]
    \centering
    \includegraphics[width=\textwidth]{2_publication_trends.pdf}
    \caption{Annual publication frequency of included studies (2016--2025), illustrating an increase in research activity from 2021 onwards.}
    \label{fig:publication-trends}
\end{figure}

Methodological approaches were diverse. Pilot studies ($n=10$, $40\%$) and prototype development studies ($n=8$, $32\%$) were the most common designs. Studies were most frequently conducted among community-dwelling older adults ($n=9$, $36\%$), followed by hospital and LTC facilities. Sample sizes varied considerably from two to 131 participants (median=28).

Based on their primary goal, the included studies were classified into four distinct application categories:

\subsection{Behaviour Change and Recommendation ($n=11$)}

Systems focused on providing personalised guidance and support to encourage healthier dietary behaviours \citep{balsaUsabilityIntelligentVirtual2020,chatterjeeDesigningInternetofThingsIoT2018,chopraDigiMateLeveragingLarge2024,chunDeepLearningBased2024,gautamFoodAssessmentModel2022,obaidAutomaticFoodintakeMonitoring2018,petrauskasExplainableArtificialIntelligencebased2021,pfistererEnhancingFoodIntake2022,rameshNovelRulebasedRecommender2021,suryawanshiMalnutritionIdentificationGeriatric2021,wangAreNutritionalLabeling2024}. Technological approaches were diverse, with deep learning being the most common ($n=4$). Most systems were in early developmental stages, classified as prototypes ($n=6$) or pilot studies ($n=4$). Only one study was an RCT, which found that an LLM-based coach increased weekly exercise and reduced fast-food consumption \citep{chopraDigiMateLeveragingLarge2024}.

\subsection{Dietary Assessment and Monitoring ($n=6$)}

Applications aimed to automate the tracking and quantification of food and nutrient intake \citep{chewADLDigitalToolkit2025,dimartinoMalnutritionRiskAssessment2021,grecoEmotionalFeaturesInteractions2021,khattakSelfAdaptiveSystem2016,kobayashiFrailtyMeterHealthy2021,seintMedicationMealIntake2018}. Systems primarily employed deep learning for image-based food recognition and general machine learning for analysing sensor and clinical data. These systems were frequently multimodal, combining inputs such as meal photos, sensor data, and text. One exception was an operationally tested system that used machine learning on clinical and dietary data in a LTC facility, achieving a median accuracy of $94\%$ and recall of $92\%$ in detecting malnutrition risk \citep{dimartinoMalnutritionRiskAssessment2021}.

\subsection{Screening and Risk Prediction ($n=4$)}

Systems designed to identify specific health risks or predict user behaviour \citep{haiSelfFeedingRobotElder2022,liKGSCSSmartCare2024,papathanailEvaluationNovelArtificial2021,spoorenImprovingRecommendationsMeals2023}. All systems were prototypes or pilot studies. The image-based dietary assessment system demonstrated a mean relative error of $11.64\%$ for energy intake estimation \citep{papathanailEvaluationNovelArtificial2021}.

\subsection{Activity Recognition ($n=4$)}

Systems focused on automated detection of fine-grained eating behaviours using wearable sensors or video \citep{ioakeimidisCanBiteDetection2024,kyritsisAssessmentRealLife2021,seokFeasibilityEfficacyTouchCare2022,wangDetectionEatingGestures2024}. The dominant approach was deep learning with emphasis on sequential data models (RNNs, LSTMs, TCNs). One study measured clinical impact, finding that an AI-driven monitoring system led to statistically significant improvement in participants' nutritional status scores over five months \citep{seokFeasibilityEfficacyTouchCare2022}.

\subsection{Cross-Cutting Data Infrastructure Themes}

Reliance on multimodal data pipelines was evident, with $68\%$ ($n=17$) of studies integrating two or more data types. Image data ($n=11$) was central for dietary assessment. Sensor data ($n=11$) was used for behaviour and physiological monitoring. Tabular data ($n=11$) served as the foundation for personalisation and clinical risk modelling. Text data ($n=7$) was central to conversational agents and knowledge-based systems. Figure~\ref{fig:modality-heatmap} visualises the distribution of these modalities across application domains; Figure~\ref{fig:technical-pipeline-sankey} traces the dominant technological pathways from input modality through AI method to application goal.

\begin{figure}[htbp]
    \centering
    \includegraphics[width=\textwidth]{3_modality_application_heatmap.pdf}
    \caption{Heatmap of data modality usage across application domains. Colour intensity corresponds to the number of studies, highlighting the dominance of tabular data in behaviour change systems and the reliance on sensor and image data for monitoring and activity recognition tasks.}
    \label{fig:modality-heatmap}
\end{figure}

\begin{figure}[htbp]
    \centering
    \includegraphics[width=\textwidth]{4_technical_pipeline_sankey.pdf}
    \caption{Multi-level Sankey diagram mapping the technical pipeline from Input Data Modality $\rightarrow$ AI Technology $\rightarrow$ Application Goal. The visualisation highlights dominant technological pathways, such as the reliance on deep learning for image-based dietary assessment and the use of machine learning with tabular data for behaviour change.}
    \label{fig:technical-pipeline-sankey}
\end{figure}

Data availability remained a significant limitation. Only four studies ($16\%$) utilised publicly available datasets, while the majority ($84\%$) relied on private, proprietary, or ``available upon request'' data, potentially hindering reproducibility.

A defining feature of the reviewed literature is its nascent nature. $92\%$ of systems were either prototypes ($n=11$) or undergoing pilot testing ($n=12$). Only two systems had advanced beyond this preliminary phase. Figure~\ref{fig:evidence-gap-map} presents the full distribution of studies across application goals, study settings, and system maturity levels.

\begin{figure}[htbp]
    \centering
    \includegraphics[width=\textwidth]{5_evidence_gap_map.pdf}
    \caption{Evidence Gap Map illustrating the distribution of studies across Application Goals and Study Settings. Bubble size represents the number of studies, and colour indicates the dominant system maturity stage (e.g., Prototype, Pilot). The map highlights a lack of research in formal care settings for several application domains.}
    \label{fig:evidence-gap-map}
\end{figure}

\section{Discussion of Review Findings}
\label{sec:discussion-review}

This scoping review of 25 studies indicates that AI in geriatric nutrition is a nascent field characterised by a gap between technological potential and real-world application. Research is concentrated on applications for behaviour change and recommendation ($n=11$) and dietary assessment and monitoring ($n=6$), accounting for $68\%$ of the included literature. The majority of systems ($n=23$, $92\%$) remain in prototype or pilot stages, with corresponding gaps in real-world validation and data availability.

The evidence gap map (Figure~\ref{fig:evidence-gap-map}) reveals opportunities and disparities. Applications for behaviour change and recommendation are predominantly studied in community-dwelling contexts, while activity recognition and screening are more common in hospital and acute care settings. However, no studies focused on behaviour change and recommendation were conducted in LTC settings---a notable void given the high prevalence of malnutrition in this demographic.

This translational challenge is compounded by the reliance on early-stage evaluations. With $92\%$ of studies describing prototypes or pilot tests, often utilising small, non-representative samples, the generalisability of current findings is limited. One usability study noted that their ``participant sample was skewed towards expert technology users'' \citep{balsaUsabilityIntelligentVirtual2020}, highlighting a need for pragmatic, multi-site RCTs.

Even within controlled evaluations, there is often a disparity between theoretical capability and practical utility. While technical performance in laboratory settings is often high (accuracies between $89\%$ and $97\%$), systems frequently struggled with environmental factors such as ``reflective food surfaces'' and the burden of obtaining ground-truth data \citep{papathanailEvaluationNovelArtificial2021}.

Compounding these environmental challenges are intrinsic barriers regarding data quality and the complexity of data collection in geriatric populations. Physical impairments and the heterogeneity of older adults often resulted in inconsistent data capture; CV models struggled to segment low-density or stacked foods \citep{pfistererEnhancingFoodIntake2022}.

A recurrent theme was the mismatch between system design and the specific needs and abilities of older adults. User feedback highlighted ``technical issues such as small buttons'' and ``content issues like repetitive questions and unclear language'' \citep{balsaUsabilityIntelligentVirtual2020}. This underscores the need to embed user-centric and co-design methodologies from the beginning.

Successful adoption also hinges on integration into broader clinical workflows. Many systems are presently developed as isolated applications with ``restricted'' scope \citep{chopraDigiMateLeveragingLarge2024}. Few studies addressed how AI-generated insights would connect with EHRs or formal care plans.

Progress is further hampered by a lack of data access. Only $16\%$ of studies used public datasets, impeding reproducibility and benchmarking.

\subsection{Methodological Limitations of the Review}

The restriction to English-language publications may have excluded relevant developments from non-Anglophone regions. The exclusion of grey literature likely omits the most recent commercial prototypes. While steps were taken to minimise bias through verification of random subsets, the primary screening was conducted by a single reviewer. Consistent with scoping review methodology, a formal critical appraisal of study quality was not performed. The heterogeneity of identified systems precluded a quantitative meta-analysis.

\subsection{Summary}

AI applications in geriatric nutrition are in an early stage of development. While technical capabilities are advancing, $92\%$ of identified systems remain in the prototype or pilot phase. Advancing the field requires addressing limitations in study design and data transparency. Future efforts should move beyond feasibility testing toward larger-scale evaluations, address the lack of interventional studies in LTC settings, and ensure interoperability with clinical workflows. Successful adoption depends on prioritising user-centred design and demonstrating quantifiable outcomes for the nutritional health of older adults.

\section{Implications for System Design}
\label{sec:implications}

The findings of this scoping review carry direct implications for the design of the artifact presented in the subsequent chapters. Rather than treating the review as an isolated academic exercise, this section traces explicit threads from the evidence synthesis to the design objectives that govern the artifact's architecture (Chapter~\ref{ch:methodology}). These threads constitute the thesis's central argumentative spine: every architectural decision in the artifact can be traced back to a gap or limitation identified in the reviewed literature.

\subsection{From Translational Gap to Deployment-Conscious Architecture}
\label{sec:gap-to-arch}

The review revealed that $92\%$ of systems remain at prototype or pilot stages (Section~\ref{sec:results}), with only two systems advancing beyond preliminary development \citep{dimartinoMalnutritionRiskAssessment2021,wangAreNutritionalLabeling2024}. This finding exposes a field-level failure to design for deployment. Systems that never leave the laboratory cannot generate the longitudinal evidence required to establish clinical validity.

This gap directly motivated the artifact's architectural decisions: a cloud-native deployment model (progressive web application on an edge platform, asynchronous API server on a managed cloud host, encrypted relational database on a managed data platform), configurable per-session clinical parameters that can be adjusted without code changes (Objective 2, Section~\ref{sec:objectives}), and a streaming event architecture that provides real-time feedback to sustain user engagement during multi-stage processing. While the artifact itself remains a research prototype --- and this tension is acknowledged honestly in Section~\ref{sec:limitations} --- its architecture is designed to \emph{permit} deployment rather than \emph{preclude} it.

\subsection{From Usability Mismatch to Geriatric-Centric Design}

A recurrent theme across the reviewed studies was the mismatch between system design and the abilities of older adults. Balsa et al.\ \citep{balsaUsabilityIntelligentVirtual2020} documented user complaints about ``small buttons'' and ``repetitive questions'' in a virtual assistant for diabetic older adults. Chew et al.\ \citep{chewADLDigitalToolkit2025} reported moderate usability scores (SUS 70.7) for a multimodal toolkit, suggesting room for improvement in interface design. Spooren et al.\ \citep{spoorenImprovingRecommendationsMeals2023} found that reducing the number of clicks required to log a meal improved adoption among older users.

These findings necessitated two specific design responses. First, the artifact's interaction model must minimise the number of required inputs while maximising information capture --- leading directly to \textbf{Objective 6} (Multimodal Input Fusion), which allows users to submit a photograph and voice description rather than navigating structured forms. Second, when the system does require user interaction (clarification), those interactions must be bounded and conversational rather than exhaustive --- leading directly to \textbf{Objective 4} (Bounded Clarification, $N_{max}=2$). The ``Tap-to-Toggle'' voice capture and haptic feedback channels (Section~\ref{sec:interface-design}) are further responses to the specific motor and sensory barriers documented by Wildenbos et al.\ \citep{wildenbosAgingBarriersInfluencing2018}.

\subsection{From Single-Shot Limitations to Agentic Clarification}

The review identified that image-based dietary assessment systems achieve high technical accuracy in controlled settings --- Papathanail et al.\ \citep{papathanailEvaluationNovelArtificial2021} reported $11.64\%$ mean relative error for energy intake estimation, and Gautam et al.\ \citep{gautamFoodAssessmentModel2022} achieved $97\%$ food classification accuracy --- but these systems uniformly rely on single-pass inference. The review's discussion (Section~\ref{sec:discussion-review}) noted that systems ``struggled with environmental factors'' and that the burden of obtaining ground-truth data was ``often prohibitive.''

The fundamental limitation is that a static image cannot resolve \emph{hidden variables}: preparation methods, occluded ingredients, and portion uncertainties that are invisible to computer vision. This limitation, consistently documented across the dietary assessment category ($n=6$), motivated the artifact's core innovation: the shift from passive, single-shot estimation to \emph{active, conditional clarification}. Rather than accepting the information loss inherent in image-only analysis, the artifact introduces a feedback loop --- the agentic clarification cycle --- that queries the user only when the missing information has clinical significance. This is operationalised through \textbf{Objective 1} (Quantifiable Complexity Assessment), which provides the mechanism to \emph{measure} ambiguity, and \textbf{Objective 5} (Semantic Before Quantitative), which ensures that the most impactful ambiguities are resolved first.

\subsection{From Absent Clinical Routing to Deterministic Threshold Gating}

No system in the reviewed literature implemented a formal mechanism for deciding \emph{when} to request clarification from the user. Conversational agents either always asked follow-up questions (e.g., \citep{balsaUsabilityIntelligentVirtual2020,grecoEmotionalFeaturesInteractions2021}) or never did (e.g., passive image analysis systems). This binary approach --- always-ask or never-ask --- is precisely the manifestation of the friction-fidelity paradox at the system design level.

The absence of a middle ground motivated the artifact's most distinctive contribution: the \textbf{Deterministic Complexity Scoring Engine} coupled with \textbf{Clinical Threshold Routing} (Objectives 1--3). By computing an auditable complexity score and comparing it against a configurable threshold, the system creates a continuous spectrum between silence and interrogation. The three-tier routing policy (mandatory override, clinical threshold, confidence gate) ensures that the system's behaviour is governed by transparent arithmetic rather than opaque heuristics. This design directly addresses the review's finding that systems lack integration with clinical workflows (Section~\ref{sec:discussion-review}) by embedding clinical sensitivity as a first-class parameter.

\subsection{From Data Opacity to Public Benchmark Evaluation}

The review found that $84\%$ of studies relied on private, proprietary, or ``available upon request'' datasets (Section~\ref{sec:results}). This data opacity impedes reproducibility and makes it impossible to compare systems against common baselines.

This finding motivated the decision to evaluate the artifact against the publicly available Nutrition5k dataset \citep{thamesNutrition5kAutomaticNutritional2021}. While the dataset has its own limitations (Western dietary bias, controlled photography conditions), its public availability ensures that the reported metrics ($22.8\%$ overall MAE reduction, $92\%$ TNR on Simple items) are independently verifiable and comparable to future work.

The implications outlined above translate directly into the design objectives formalised in the following chapter. Where this chapter asked \emph{what does the field need?}, Chapter~\ref{ch:methodology} asks \emph{what must the artifact do?} --- and answers with six verifiable requirements derived from the evidence gaps documented here.
